\documentclass[10pt,a4paper]{amsart}
\usepackage{amsmath,amssymb,amsthm,graphicx}
% Retained minimal preamble for the Stacks Project chapter more-morphisms.tex (More on Morphisms).
% Verbatim excerpt of the chapter preamble: Xy-pic setup (native preamble lines 14-18) and the native
% theorem declarations with the mathematical macros (native preamble lines 154-217). Copied mechanically;
% nothing added, removed or redefined. The wrapper class and the amsmath/amssymb/amsthm/graphicx base
% packages are supplied by the assembler, so the class/externaldocument part of the native preamble is omitted.
% For commutative diagrams we use Xy-pic
\usepackage[all]{xy}

% We use 2cell for 2-commutative diagrams.
\xyoption{2cell}
% Theorem environments.
%
\theoremstyle{plain}
\newtheorem{theorem}[subsection]{Theorem}
\newtheorem{proposition}[subsection]{Proposition}
\newtheorem{lemma}[subsection]{Lemma}

\theoremstyle{definition}
\newtheorem{definition}[subsection]{Definition}
\newtheorem{example}[subsection]{Example}
\newtheorem{exercise}[subsection]{Exercise}
\newtheorem{situation}[subsection]{Situation}

\theoremstyle{remark}
\newtheorem{remark}[subsection]{Remark}
\newtheorem{remarks}[subsection]{Remarks}

\numberwithin{equation}{subsection}

% Macros
%
\def\lim{\mathop{\mathrm{lim}}\nolimits}
\def\colim{\mathop{\mathrm{colim}}\nolimits}
\def\Spec{\mathop{\mathrm{Spec}}}
\def\Hom{\mathop{\mathrm{Hom}}\nolimits}
\def\Ext{\mathop{\mathrm{Ext}}\nolimits}
\def\SheafHom{\mathop{\mathcal{H}\!\mathit{om}}\nolimits}
\def\SheafExt{\mathop{\mathcal{E}\!\mathit{xt}}\nolimits}
\def\Sch{\mathit{Sch}}
\def\Mor{\mathop{\mathrm{Mor}}\nolimits}
\def\Ob{\mathop{\mathrm{Ob}}\nolimits}
\def\Sh{\mathop{\mathit{Sh}}\nolimits}
\def\NL{\mathop{N\!L}\nolimits}
\def\CH{\mathop{\mathrm{CH}}\nolimits}
\def\proetale{{pro\text{-}\acute{e}tale}}
\def\etale{{\acute{e}tale}}
\def\QCoh{\mathit{QCoh}}
\def\Ker{\mathop{\mathrm{Ker}}}
\def\Im{\mathop{\mathrm{Im}}}
\def\Coker{\mathop{\mathrm{Coker}}}
\def\Coim{\mathop{\mathrm{Coim}}}

% Boxtimes
%
\DeclareMathSymbol{\boxtimes}{\mathbin}{AMSa}{"02}

%
% Macros for moduli stacks/spaces
%
\def\QCohstack{\mathcal{QC}\!\mathit{oh}}
\def\Cohstack{\mathcal{C}\!\mathit{oh}}
\def\Spacesstack{\mathcal{S}\!\mathit{paces}}
\def\Quotfunctor{\mathrm{Quot}}
\def\Hilbfunctor{\mathrm{Hilb}}
\def\Curvesstack{\mathcal{C}\!\mathit{urves}}
\def\Polarizedstack{\mathcal{P}\!\mathit{olarized}}
\def\Complexesstack{\mathcal{C}\!\mathit{omplexes}}
% \Pic is the operator that assigns to X its picard group, usage \Pic(X)
% \Picardstack_{X/B} denotes the Picard stack of X over B
% \Picardfunctor_{X/B} denotes the Picard functor of X over B
\def\Pic{\mathop{\mathrm{Pic}}\nolimits}
\def\Picardstack{\mathcal{P}\!\mathit{ic}}
\def\Picardfunctor{\mathrm{Pic}}
\def\Deformationcategory{\mathcal{D}\!\mathit{ef}}
\usepackage{cleveref}
\crefname{originalreference}{original reference}{original references}
\Crefname{originalreference}{Original reference}{Original references}
\makeatletter
\newlabel{algebra-lemma-normal-goes-up}{{033C}{source}{}{sourcepointer-1}{}}
\newlabel{algebra-lemma-normal-goes-up@cref}{{[originalreference][0][]033C}{[1][0][]source}{}{}{}}
\makeatother
\crefname{originalreference}{original reference}{original references}
\Crefname{originalreference}{Original reference}{Original references}
\makeatletter
\newlabel{algebra-lemma-separable-smooth}{{00TV}{source}{}{sourcepointer-2}{}}
\newlabel{algebra-lemma-separable-smooth@cref}{{[originalreference][0][]00TV}{[1][0][]source}{}{}{}}
\makeatother
\crefname{originalequation}{original equation}{original equations}
\Crefname{originalequation}{Original equation}{Original equations}
\makeatletter
\newlabel{equation-normalized-base-change}{{09IK}{source}{}{sourcepointer-3}{}}
\newlabel{equation-normalized-base-change@cref}{{[originalequation][0][]09IK}{[1][0][]source}{}{}{}}
\makeatother
\crefname{originallemma}{original lemma}{original lemmas}
\Crefname{originallemma}{Original lemma}{Original lemmas}
\makeatletter
\newlabel{lemma-nowhere-dense-generic-fibre}{{054X}{source}{}{sourcepointer-4}{}}
\newlabel{lemma-nowhere-dense-generic-fibre@cref}{{[originallemma][0][]054X}{[1][0][]source}{}{}{}}
\makeatother
\crefname{originalreference}{original reference}{original references}
\Crefname{originalreference}{Original reference}{Original references}
\makeatletter
\newlabel{more-algebra-lemma-epp-essentially-finite-type-separable}{{0BRP}{source}{}{sourcepointer-5}{}}
\newlabel{more-algebra-lemma-epp-essentially-finite-type-separable@cref}{{[originalreference][0][]0BRP}{[1][0][]source}{}{}{}}
\makeatother
\crefname{originalreference}{original reference}{original references}
\Crefname{originalreference}{Original reference}{Original references}
\makeatletter
\newlabel{more-algebra-lemma-solution-goes-up}{{0GLR}{source}{}{sourcepointer-6}{}}
\newlabel{more-algebra-lemma-solution-goes-up@cref}{{[originalreference][0][]0GLR}{[1][0][]source}{}{}{}}
\makeatother
\crefname{originalreference}{original reference}{original references}
\Crefname{originalreference}{Original reference}{Original references}
\makeatletter
\newlabel{more-algebra-proposition-epp-essentially-finite-type}{{09II}{source}{}{sourcepointer-7}{}}
\newlabel{more-algebra-proposition-epp-essentially-finite-type@cref}{{[originalreference][0][]09II}{[1][0][]source}{}{}{}}
\makeatother
\crefname{originalreference}{original reference}{original references}
\Crefname{originalreference}{Original reference}{Original references}
\makeatletter
\newlabel{morphisms-definition-smooth}{{01V5}{source}{}{sourcepointer-8}{}}
\newlabel{morphisms-definition-smooth@cref}{{[originalreference][0][]01V5}{[1][0][]source}{}{}{}}
\makeatother
\crefname{originalreference}{original reference}{original references}
\Crefname{originalreference}{Original reference}{Original references}
\makeatletter
\newlabel{morphisms-lemma-set-points-where-fibres-smooth}{{02V4}{source}{}{sourcepointer-9}{}}
\newlabel{morphisms-lemma-set-points-where-fibres-smooth@cref}{{[originalreference][0][]02V4}{[1][0][]source}{}{}{}}
\makeatother
\crefname{originaltheorem}{original theorem}{original theorems}
\Crefname{originaltheorem}{Original theorem}{Original theorems}
\makeatletter
\newlabel{theorem-normalized-base-change-with-reduced-fibre}{{09IL}{source}{}{sourcepointer-10}{}}
\newlabel{theorem-normalized-base-change-with-reduced-fibre@cref}{{[originaltheorem][0][]09IL}{[1][0][]source}{}{}{}}
\makeatother
\crefname{originalreference}{original reference}{original references}
\Crefname{originalreference}{Original reference}{Original references}
\makeatletter
\newlabel{varieties-lemma-generic-points-geometrically-reduced}{{04KS}{source}{}{sourcepointer-11}{}}
\newlabel{varieties-lemma-generic-points-geometrically-reduced@cref}{{[originalreference][0][]04KS}{[1][0][]source}{}{}{}}
\makeatother
\crefname{originalreference}{original reference}{original references}
\Crefname{originalreference}{Original reference}{Original references}
\makeatletter
\newlabel{varieties-lemma-geometrically-reduced-dense-smooth-open}{{056V}{source}{}{sourcepointer-12}{}}
\newlabel{varieties-lemma-geometrically-reduced-dense-smooth-open@cref}{{[originalreference][0][]056V}{[1][0][]source}{}{}{}}
\makeatother
\title{Normalized base change along a separable residue extension is smooth at generic fibre points}
\author{The Stacks Project Authors}
\date{stacks-project:more-morphisms.tex:local-snapshot-20261011; GNU FDL 1.2 or later}
\begin{document}
\maketitle
\noindent\textit{Editorial note.} Counted span is the author's own proof environment, source lines 19412-19478, immediately following the statement at 19400-19410.\par\medskip
\noindent\textit{Editorial note.} The proof opens with 'This is proved in exactly the same manner as [the general theorem]' and then reproduces the argument for the separable case; the reproduced argument is inside the span.\par\medskip
\noindent\textit{Editorial note.} Nine cross-chapter pointers (More on Algebra, Algebra, Morphisms, Varieties, Decent Spaces) are recorded in assembly.external\_references as navigation pointers.\par\medskip
\noindent\textit{Editorial note.} Steps the author leaves to the reader are flagged here for the reviewer: source line(s) 19467 carry an "omitted"/"follows easily"/"immediate from the definitions" phrase. Those sentences are retained verbatim; the displayed argument around them stays inside the counted span.\par\medskip
\noindent\textit{Allowed standard prerequisites (editorial).}\begin{itemize}
\item Normalized base change of the general theorem. chapter theorem theorem-normalized-base-change-with-reduced-fibre, cited at proof line 19412 onward. Supplies the framework (normalization T, normalized base change Y) whose smoothness locus is refined here.
\item Separable field extensions are smooth. Algebra, lemma algebra-lemma-separable-smooth, cited at proof line 19476. Converts separability of the residue field extension at a point into smoothness of the fibre there.
\item Solutions of an essentially finite type problem go up. More on Algebra, lemmas more-algebra-lemma-epp-essentially-finite-type-separable, more-algebra-lemma-solution-goes-up and proposition more-algebra-proposition-epp-essentially-finite-type, cited throughout the proof. Provides the algebra result that a separable solution to the essentially finite type problem exists after a finite separable extension.
\item Where the fibres are smooth and generically reduced. Morphisms, lemma morphisms-lemma-set-points-where-fibres-smooth and Varieties, lemma varieties-lemma-generic-points-geometrically-reduced, cited in the proof. Identifies the loci of points where the fibres are smooth and controls generic points of the fibres.
\end{itemize}\medskip
\ifdefined\hypertarget\hypertarget{sourcepointer-1}{}\fi
\noindent\textit{Source reference.} Original reference 033C: focused-seeds/geometry/S9-morph-normalized-base-change-separable/ref-tags.txt, source lines 4--4. Navigation pointer only: the counted proof cites this statement by \textbackslash{}ref as a standing input; the deduction the counted proof has to make is carried out inside the counted span, and no part of that deduction is delegated to this pointer.\par\medskip
\ifdefined\hypertarget\hypertarget{sourcepointer-2}{}\fi
\noindent\textit{Source reference.} Original reference 00TV: focused-seeds/geometry/S9-morph-normalized-base-change-separable/ref-tags.txt, source lines 5--5. Navigation pointer only: the counted proof cites this statement by \textbackslash{}ref as a standing input; the deduction the counted proof has to make is carried out inside the counted span, and no part of that deduction is delegated to this pointer.\par\medskip
\ifdefined\hypertarget\hypertarget{sourcepointer-3}{}\fi
\noindent\textit{Source reference.} Original equation 09IK: sources/stacks-more-morphisms/pinned-original-source.tex, source lines 19260--19260. Navigation pointer only: the counted proof cites this statement by \textbackslash{}ref as a standing input; the deduction the counted proof has to make is carried out inside the counted span, and no part of that deduction is delegated to this pointer.\par\medskip
\ifdefined\hypertarget\hypertarget{sourcepointer-4}{}\fi
\noindent\textit{Source reference.} Original lemma 054X: sources/stacks-more-morphisms/pinned-original-source.tex, source lines 6524--6524. Navigation pointer only: the counted proof cites this statement by \textbackslash{}ref as a standing input; the deduction the counted proof has to make is carried out inside the counted span, and no part of that deduction is delegated to this pointer.\par\medskip
\ifdefined\hypertarget\hypertarget{sourcepointer-5}{}\fi
\noindent\textit{Source reference.} Original reference 0BRP: focused-seeds/geometry/S9-morph-normalized-base-change-separable/ref-tags.txt, source lines 6--6. Navigation pointer only: the counted proof cites this statement by \textbackslash{}ref as a standing input; the deduction the counted proof has to make is carried out inside the counted span, and no part of that deduction is delegated to this pointer.\par\medskip
\ifdefined\hypertarget\hypertarget{sourcepointer-6}{}\fi
\noindent\textit{Source reference.} Original reference 0GLR: focused-seeds/geometry/S9-morph-normalized-base-change-separable/ref-tags.txt, source lines 7--7. Navigation pointer only: the counted proof cites this statement by \textbackslash{}ref as a standing input; the deduction the counted proof has to make is carried out inside the counted span, and no part of that deduction is delegated to this pointer.\par\medskip
\ifdefined\hypertarget\hypertarget{sourcepointer-7}{}\fi
\noindent\textit{Source reference.} Original reference 09II: focused-seeds/geometry/S9-morph-normalized-base-change-separable/ref-tags.txt, source lines 8--8. Navigation pointer only: the counted proof cites this statement by \textbackslash{}ref as a standing input; the deduction the counted proof has to make is carried out inside the counted span, and no part of that deduction is delegated to this pointer.\par\medskip
\ifdefined\hypertarget\hypertarget{sourcepointer-8}{}\fi
\noindent\textit{Source reference.} Original reference 01V5: focused-seeds/geometry/S9-morph-normalized-base-change-separable/ref-tags.txt, source lines 9--9. Navigation pointer only: the counted proof cites this statement by \textbackslash{}ref as a standing input; the deduction the counted proof has to make is carried out inside the counted span, and no part of that deduction is delegated to this pointer.\par\medskip
\ifdefined\hypertarget\hypertarget{sourcepointer-9}{}\fi
\noindent\textit{Source reference.} Original reference 02V4: focused-seeds/geometry/S9-morph-normalized-base-change-separable/ref-tags.txt, source lines 10--10. Navigation pointer only: the counted proof cites this statement by \textbackslash{}ref as a standing input; the deduction the counted proof has to make is carried out inside the counted span, and no part of that deduction is delegated to this pointer.\par\medskip
\ifdefined\hypertarget\hypertarget{sourcepointer-10}{}\fi
\noindent\textit{Source reference.} Original theorem 09IL: sources/stacks-more-morphisms/pinned-original-source.tex, source lines 19282--19282. Navigation pointer only: the counted proof cites this statement by \textbackslash{}ref as a standing input; the deduction the counted proof has to make is carried out inside the counted span, and no part of that deduction is delegated to this pointer.\par\medskip
\ifdefined\hypertarget\hypertarget{sourcepointer-11}{}\fi
\noindent\textit{Source reference.} Original reference 04KS: focused-seeds/geometry/S9-morph-normalized-base-change-separable/ref-tags.txt, source lines 11--11. Navigation pointer only: the counted proof cites this statement by \textbackslash{}ref as a standing input; the deduction the counted proof has to make is carried out inside the counted span, and no part of that deduction is delegated to this pointer.\par\medskip
\ifdefined\hypertarget\hypertarget{sourcepointer-12}{}\fi
\noindent\textit{Source reference.} Original reference 056V: focused-seeds/geometry/S9-morph-normalized-base-change-separable/ref-tags.txt, source lines 12--12. Navigation pointer only: the counted proof cites this statement by \textbackslash{}ref as a standing input; the deduction the counted proof has to make is carried out inside the counted span, and no part of that deduction is delegated to this pointer.\par\medskip
\begin{lemma}[Variant with separable extension]
\label{lemma-normalized-base-change-with-reduced-fibre-separable}
Let $A$ be a Dedekind ring with fraction field $K$.
Let $X$ be a scheme flat and of finite type over $A$.
Assume $A$ is a Nagata ring and that for every generic point
$\eta$ of an irreducible component of $X$ the field
extension $\kappa(\eta)/K$ is separable.
Then there exists a finite separable extension $L/K$ such that
the normalized base change $Y$ is smooth over $\Spec(B)$
at all generic points of all fibres.
\end{lemma}

\begin{proof}
This is proved in exactly the same manner as
Theorem \ref{theorem-normalized-base-change-with-reduced-fibre}
with a few minor modifications. The most important change
is to use More on Algebra, Lemma
\ref{more-algebra-lemma-epp-essentially-finite-type-separable}
instead of More on Algebra, Proposition
\ref{more-algebra-proposition-epp-essentially-finite-type}.
During the proof we will repeatedly use that formation of the set of points
where a (flat, finitely presented) morphism like $X \to \Spec(A)$ is
smooth commutes with base change, see
Morphisms, Lemma \ref{morphisms-lemma-set-points-where-fibres-smooth}.

\medskip\noindent
Since $X$ is flat over $A$ every generic point $\eta$ of $X$ maps to the
generic point of $\Spec(A)$.
After replacing $X$ by its reduction we may assume $X$ is reduced.
In this case $X_K$ is geometrically reduced over $K$
by Varieties, Lemma \ref{varieties-lemma-generic-points-geometrically-reduced}.
Hence $X_K \to \Spec(K)$ is smooth on a dense open by
Varieties, Lemma \ref{varieties-lemma-geometrically-reduced-dense-smooth-open}.
Thus the smooth locus $U \subset X$ of the morphism $X \to \Spec(A)$
is open (by Morphisms, Definition \ref{morphisms-definition-smooth})
and is dense in the generic fibre. This reduces us to the situation
of the following paragraph.

\medskip\noindent
Assume $X$ is normal and the smooth locus $U \subset X$ of $X \to \Spec(A)$
is dense in the generic fibre. This implies that $U$ is dense in all but
finitely many fibres, see Lemma \ref{lemma-nowhere-dense-generic-fibre}.
Let $x_1, \ldots, x_r \in X \setminus U$ be the finitely many generic
points of irreducible components of $X \setminus U$ which are moreover
generic points of irreducible components of fibres of $X \to \Spec(A)$.
Set $\mathcal{O}_i = \mathcal{O}_{X, x_i}$. Observe that the fraction
field of $\mathcal{O}_i$ is the residue field of a generic point of $X$.
Let $A_i$ be the localization of $A$ at the maximal ideal corresponding
to the image of $x_i$ in $\Spec(A)$. We may apply More on Algebra, Lemma
\ref{more-algebra-lemma-epp-essentially-finite-type-separable}
and we find finite separable extensions $K_i/K$ which are
solutions for $A_i \to \mathcal{O}_i$. Let $L/K$ be a finite
separable extension dominating all of the extensions $K_i/K$.
Then $L/K$ is still a solution for $A_i \to \mathcal{O}_i$ by
More on Algebra, Lemma \ref{more-algebra-lemma-solution-goes-up}.

\medskip\noindent
Consider the diagram (\ref{equation-normalized-base-change})
with the extension $L/K$ we just produced. Note that $U_B \subset X_B$
is smooth over $B$, hence normal (for example use
Algebra, Lemma \ref{algebra-lemma-normal-goes-up}).
Thus $Y \to X_B$ is an isomorphism over $U_B$.
Let $y \in Y$ be a generic point of an irreducible
component of a fibre of $Y \to \Spec(B)$ lying over the maximal
ideal $\mathfrak m \subset B$. Assume that $y \not \in U_B$.
Then $y$ maps to one of the points $x_i$. It follows that
$\mathcal{O}_{Y, y}$ is a local ring of the integral closure
of $\mathcal{O}_i$ in $R(X) \otimes_K L$ (details omitted).
Hence because $L/K$ is a solution for
$A_i \to \mathcal{O}_i$ we see that
$B_\mathfrak m \to \mathcal{O}_{Y, y}$ is formally smooth
(this is the definition of being a "solution").
In other words, $\mathfrak m\mathcal{O}_{Y, y} = \mathfrak m_y$
and the residue field extension is separable. Hence the local ring
of the fibre at $y$ is $\kappa(y)$.
This implies the fibre is smooth over $\kappa(\mathfrak m)$
at $y$ for example by Algebra, Lemma \ref{algebra-lemma-separable-smooth}.
This finishes the proof.
\end{proof}

\end{document}
